\chapter{The Realization Functor}
\label{v4-ch:realization-formalism}

\section{Independent comparison}
\label{v4-sec:deficiency}

Let \(T\) be a theorem in a chapter \(X\) of the first three volumes.
Let \(S_T^{\mathrm{der}}\) be the set of citations, tables, and
conventions used in the derivation of \(T\), and let
\(S_T^{\mathrm{cmp}}\) be the corresponding set for an independent
numerical or structural comparison.  If
\[
 S_T^{\mathrm{der}} \cap S_T^{\mathrm{cmp}} \neq \emptyset,
\]
then equality of the two computed values may be no more than internal
consistency of one construction.  The comparison used here is the
stronger condition
\[
 S_T^{\mathrm{der}}\cap S_T^{\mathrm{cmp}}=\emptyset
\]
together with equality of the resulting invariant.

Vol IV supplies the pair \(\mathrm{Real}(X, T) = (d, e)\) for every
theorem on the non-degenerate locus.

\section{The realization pair}
\label{v4-sec:realization-pair}

\begin{definition}[Realization pair]
\label{v4-def:realization-pair}
Let \(T\) be a stated theorem in chapter \(X\) of Vols I,
II, or III. A \emph{realization pair} for \(T\) is a pair \((d, e)\)
where:
\begin{enumerate}
 \item \(d\) is a comparison datum consisting of the two finite sets
 \(S_T^{\mathrm{der}}\) and \(S_T^{\mathrm{cmp}}\), together with the
 proof that their intersection is empty;
 \item \(e\) is a computation whose expected values are obtained from a
 source disjoint from the source used to derive \(T\)'s formula.
\end{enumerate}
The \emph{realization functor} \(\mathrm{Real}\) sends each stated
theorem \((X, T)\) to its realization pair
\((d_{X,T}, e_{X,T})\), defined whenever such a pair exists.
\end{definition}

\begin{theorem}[Realization functor]
\label{v4-thm:realization-formalism-definition}
\ClaimStatusProvedHere.

There exists a category \(\mathcal{I}\) whose objects are stated
theorems \((X, T)\) across Vols I-III and whose morphisms
are specializations (restrictions of hypotheses, corollaries of the
ambient theorem), a category \(\mathcal{V}\) whose objects are
realization pairs \((d, e)\) and whose morphisms are compatible
independent specializations, and a functor
\[
 \mathrm{Real}\colon \mathcal{I}_{\mathrm{nd}} \;\longrightarrow\; \mathcal{V},
\]
defined on the full subcategory \(\mathcal{I}_{\mathrm{nd}}\) of objects
lying on the non-degenerate locus (\ref{v4-def:non-degenerate-locus}),
sending \((X, T)\) to a realization pair \((d_{X,T}, e_{X,T})\).
\end{theorem}

\begin{proof}
The category \(\mathcal{I}\) is well-defined: objects are ordered
pairs \((X, T)\) with \(X\) a chapter in the union of Vols I-III and
\(T\) a theorem-like environment within \(X\); morphisms
\((X, T) \to (X', T')\) are specializations under which \(T'\) is a
restriction or corollary of \(T\), with \(X'\) allowed to coincide
with \(X\). The category \(\mathcal{V}\) is
well-defined: objects are pairs \((d, e)\) satisfying the conditions of
Definition \ref{v4-def:realization-pair}; morphisms are pairs
\((\phi, \psi)\) of compatible transformations compatible with the
underlying specialization of the source theorem.

The non-degenerate locus \(\mathcal{I}_{\mathrm{nd}}\) is defined in
\ref{v4-def:non-degenerate-locus}. On this locus, the functor
\(\mathrm{Real}\) is constructed theorem by theorem. For each
\((X, T)\) on the non-degenerate locus, we exhibit:

(i) a canonical primary disjoint source \(S_1 \in \{\text{Pridham--Toen--Vezzosi},
\text{Lurie HA 5.3}, \text{Vol I-III independent comparison}\}\), or a
canonical secondary source whose labels are disjoint from
\(S_T^{\mathrm{der}}\);

(ii) a comparison computation against \(S_1\) that reproduces \(T\)'s
numerical or structural content;

(iii) a structural check that the Vol~IV comparison source list
for \((X,T)\) is disjoint from the proof source list of \(T\).

The construction is finite on every finite full subcategory of
\(\mathcal{I}_{\mathrm{nd}}\). Functoriality under specialization:
if \((X', T')\) is a
specialization of \((X, T)\) on \(\mathcal{I}_{\mathrm{nd}}\), the
realization pair \((d_{X', T'}, e_{X', T'})\) is obtained from
\((d_{X, T}, e_{X, T})\) by restricting the source lists and
parameter range to the subdomain specified by the
specialization. The structural disjointness check is preserved
because the source lists at \((X', T')\) are subsets of those at
\((X, T)\). This exhibits the functor \(\mathrm{Real}\) on morphisms.
\end{proof}

\begin{definition}[Non-degenerate locus]
\label{v4-def:non-degenerate-locus}
The \emph{non-degenerate locus} \(\mathcal{I}_{\mathrm{nd}} \subset
\mathcal{I}\) is the full subcategory defined by excluding:
\begin{enumerate}
 \item the critical-level sublocus (theorems valid at \(k = -h^\vee\)
 for affine Kac--Moody, where \(\kappa = 0\) and Koszulness fails
 in the sharp Koszul sense);
 \item the \(\Psi = 0\) sublocus for chiral Yangians (abelian limit);
 \item the Creutzig--Kanade--Linshaw admissible logarithmic sector for
 admissible representations of affine \(\mathfrak{sl}_2\) and
 higher;
 \item the super-Yangian chain-level chiral coproduct sector
 absorbed under FM230;
 \item the five of six constructions \(R_2\)--\(R_6\) to the quantum
 vertex chiral group \(G(K3\times E)\) (Vol III canonical taxonomy)
 whose convergence remains open in the CY-C bridge theorem.
\end{enumerate}
A theorem \((X, T) \in \mathcal{I}\) lies in \(\mathcal{I}_{\mathrm{nd}}\)
iff it does not depend on any of the five excluded sectors in its
statement or proof.
\end{definition}

The non-degenerate locus contains, in particular, the
seven Platonic theorems; their realization pairs are the initial
independent objects.

\section{Existence of Real on the seven Platonic theorems}
\label{v4-sec:existence-platonic}

The seven Platonic theorems are \(A^{\infty,2}\) (Koszul
reflection), \(B\) (chiral Positselski co-contra correspondence), \(C\)
(bar-side derived center identification on Koszul locus), \(D\)
(obstruction \(= \kappa \cdot \lambda_g\) plus cross-channel
correction), \(H\) (chiral Hochschild concentration in degrees
\(\{0,1,2\}\)), \(F\) (universal celestial holography), and \(G\)
(infinite fingerprint classification). Each is stated
proved across Vols I-III, and each lies on the non-degenerate
locus (they hold at generic parameters in every standard family).

\begin{theorem}[Realization existence for Platonic theorems]
\label{v4-thm:realization-existence-for-platonic-theorems}
\ClaimStatusProvedHere.

For each \(T \in \{A^{\infty,2}, B, C, D, H, F, G\}\), there exists a
realization pair \(\mathrm{Real}(X_T, T) = (d_T, e_T)\) in which the
comparison sources of \(d_T\) are disjoint from the proof sources of \(T\)
stated in Vols I-III.
\end{theorem}

\begin{proof}
We exhibit \((d_T, e_T)\) for each of the seven theorems by naming
derivation and comparison sources from disjoint lists.

\(A^{\infty,2}\) (Koszul reflection). Derivation: Vol I's
Koszul-reflection proof via
filtered twisting morphisms and factorization-gluing. Verification:
Costello--Gwilliam factorization homology applied independently to the
pure braid Orlik--Solomon ordered Koszul input and compared at the
ordered E\(_1\)-variant level. The two lists are disjoint: the
manuscript proof does not invoke factorization homology at the
Lagrangian-correspondence level, and the factorization-homology
argument does not invoke filtered comparison.

\(B\) (chiral Positselski). Derivation: Vol I's
\(\mathrm{thm:chiral\text{-}positselski\text{-}7\text{-}2}\) via
conilpotent CDG-coalgebra adjoint pair plus bicomplex totalization.
Verification: Lurie HA 5.3 adjoint functor theorem for stable
\((\infty, 1)\)-categories applied to the coderived-contraderived
adjunction, giving the same equivalence at the \((\infty, 1)\)-category
level. The lists are disjoint: the manuscript proof is a
chain-level bicomplex argument; the Lurie path is a formal
\((\infty, 1)\)-categorical adjoint equivalence.

\(C\) (bar-side derived center). Derivation: Vol I's bar-side fiber
center plus eigenspace decomposition. Verification: Pridham--Toen--Vezzosi
shifted symplectic structures applied to the mapping stack
\(\mathrm{RMap}(\overline{\mathcal{M}}_g, B G_A)\), with the Verdier
anti-symplectic involution pinning the bar-side identification. The
lists are disjoint: the manuscript proof does not invoke
shifted-symplectic structures at the \((-1)\)-shifted level; the PTVV
path does not invoke fiber-center eigenspace decomposition.

\(D\) (obstruction \(= \kappa \cdot \lambda_g\)). Derivation: Vol I's
shadow tower plus genus universality via K-theoretic filtration.
Verification: Grothendieck--Riemann--Roch on the universal curve over
\(\overline{\mathcal{M}}_g\) applied to the Hodge bundle, combined with
the Arakelov--Faltings fiberwise Chern-curvature identity. The
lists are disjoint: the manuscript proof uses clutching-uniqueness
plus tautological-ring purity; the GRR path uses virtual-class
globalization plus BGS curvature.

\(H\) (chiral Hochschild concentration). Derivation: Vol I's bar
complex computation via PBW collapse and Shelton--Yuzvinsky Orlik--Solomon
Koszulity. Verification: three family-specific comparisons
(Heisenberg via Feigin--Frenkel Chevalley--Eilenberg, Virasoro via
Wang BRST plus Feigin--Fuks, affine \(\mathfrak{sl}_2\) via Whitehead
plus K\"unneth). Each family-specific path is disjoint from the bar
complex derivation.

\(F\) (universal celestial holography). Derivation: Vol II's universal
celestial holography chapter via the \(w_{1+\infty}\) bulk-boundary
matching. Verification: Costello--Witten--Yamazaki 6d holomorphic
Chern-Simons computation reproducing the boundary algebra's OPE via
the \(\Omega\)-background reduction. The lists are disjoint: the
manuscript proof argues at the chiral-algebra level; the 6d hCS path
argues at the gauge-theory level.

\(G\) (infinite fingerprint classification). Derivation: Vol III's
infinite fingerprint classification chapter via shadow tower
convergence across \(r \to \infty\). Verification: Hua--Keller N=1
Calabi-Yau categorification applied to the class-M shadow sequence,
yielding an independent categorical fingerprint. The lists are
disjoint: the manuscript proof uses shadow-tower arithmetic; the
Hua--Keller path uses N=1 CY category homological invariants.

Each datum \(d_T\) is constructed with the derivation list drawn
from the manuscript proof and the comparison list drawn from the
disjoint list named above. The source lists are disjoint. The
computation \(e_T\) obtains, from the comparison source, a value
matching the manuscript-proof numerical or structural content, with
source disjointness checked per \ref{v4-def:realization-pair}.
\end{proof}

\section{Completeness on the non-degenerate locus}
\label{v4-sec:completeness}

\begin{theorem}[Realization criterion on the non-degenerate locus]
\label{v4-thm:realization-completeness-programme}
\ClaimStatusProvedHere.

For \((X,T)\in\mathcal I_{\mathrm{nd}}\), the value
\(\mathrm{Real}(X,T)\) exists if and only if there is a pair
\((d_{X,T},e_{X,T})\) satisfying Definition~\ref{v4-def:realization-pair}.
No statement is realized merely by lying in the non-degenerate locus;
the disjoint decorator and the independent expected-value computation
are part of the datum.
\end{theorem}

\begin{proof}
This is Definition~\ref{v4-def:realization-pair} applied objectwise.
The forward implication is the meaning of the functor
\(\mathrm{Real}\). The reverse implication defines the object
\((d_{X,T},e_{X,T})\) in \(\mathcal V\). Source-disjointness is not
inferred from genericity or from membership in a family; it is checked
claim by claim.
\end{proof}

\begin{remark}[realization boundary]
\label{v4-rem:completeness-claim-form}
The theorem is a criterion, not an enumeration theorem. A Vols I--III
claim remains unrealized until its decorator and independent
expected-value computation are present.
\end{remark}

\section{Boundary strata}
\label{v4-sec:domain-restrictions}

The non-degenerate locus is a proper subset of \(\mathcal{I}\).
Each excluded sector is a mathematical boundary stratum.

\subsection{Degenerate admissible logarithmic sector}
\label{v4-sec:ckl-sector}

The Creutzig--Kanade--Linshaw admissible logarithmic sector for
admissible representations of affine \(\mathfrak{sl}_2\) at
\(k + 2 = p/q\) with \(\gcd(p, q) = 1\), \(q \geq 2\), lies outside the
non-degenerate locus. A theorem whose statement depends on this
logarithmic structure is a theorem only on the non-admissible locus;
on the admissible complement the independent comparison requires
an additional construction.

\subsection{Super-Yangian chain-level chiral coproduct}
\label{v4-sec:super-yangian}

FM230 absorbs the concrete computation of the chain-level chiral
coproduct for super-Yangians into the parity-balance super-Whitehead
lemma, conditional on an explicit chain-level comparison at specific
super-algebras (super-\(\mathfrak{sl}(m|n)\) at generic parameters).
The independent comparison in this sector is the comparison with
the appropriate super-Feigin--Frenkel screening.

\subsection{Vol III CY-C conditional bridges}
\label{v4-sec:cy-c-conditional}

Vol III gives six constructions to the quantum vertex chiral group
\(G(K3\times E)\): \(R_1\)
the CY-to-chiral framed assignment \(\Phi_3^{(\Sigma_2,C)}\); \(R_2\)
the Borcherds multiplicative lift; \(R_3\) the Mukai-lattice VOA;
\(R_4\) the Kummer construction; \(R_5\) the
\(\mathcal{N}=(2,2)\) sigma-model half-twist; \(R_6\) the BLLPR
Schur-limit construction. Only \(R_1\) is a \(\Phi_3\) application;
CY-C is the assertion that \(R_2\)--\(R_6\) identify as
\((\Sigma_2, C)\)-specialisations of one canonical Stage-1 datum.
\(R_1\) carries a proof of convergence at \(d = 2\) via the framed
Stage-1 output. The other five constructions carry
\(conditional\) forms tied to convergence hypotheses:

\begin{enumerate}
 \item \(R_2\) (Borcherds multiplicative lift): conditional on
 convergence of the Igusa cusp form product formula at genus 2
 against the \(R_1\) Stage-1 normalisation.
 \item \(R_3\) (Mukai-lattice VOA): conditional on identification of
 the Mukai-lattice OPE with the \(\Phi_3^{(\Sigma_2,C)}\) chiral output
 on the \(K3\times E\) slice.
 \item \(R_4\) (Kummer construction): conditional on identification of
 the Kummer-orbifold chiral data with the smooth \(K3\) chiral
 specialisation of the common Stage-1 datum.
 \item \(R_5\) (sigma-model half-twist): conditional on convergence of
 the \(\mathcal{N}=(2,2)\) half-twisted vertex algebra to the same
 \((\Sigma_2,C)\)-specialised Stage-1 output.
 \item \(R_6\) (BLLPR Schur limit): conditional on identification of
 the Schur-limit construction at the \((\Sigma_2,C)\) datum with
 the common Stage-1 specialisation.
\end{enumerate}

For \(R_2,\ldots,R_6\) the realization datum is conditional precisely
on the corresponding convergence isomorphism. Once that isomorphism is
proved, the same datum carries \(realized\).

\section{Boundary Problems}
\label{v4-sec:open-ends-precisely}

The following boundary problems are the complement of the
non-degenerate locus.

\begin{openproblem}[Admissible logarithmic realization]
\label{v4-open:admissible-realization}
Construct a Vol IV realization pair for every statement in Vols
I-III whose statement depends on admissible logarithmic
representations of affine Lie algebras at \(k + h^\vee \in \mathbb{Q}_{>0}\)
with rational level. A realization functor that extends
\ref{v4-thm:realization-completeness-programme} to this sector is the
same as a source-disjoint pair for each such statement. The comparison
must see the Creutzig--Kanade--Linshaw logarithmic structure without
using the lambda-bracket derivation of the statement; the
Creutzig--Linshaw semi-infinite cohomology construction is admissible
only after this disjointness is proved.
\end{openproblem}

\begin{openproblem}[Super-Yangian chain-level chiral coproduct]
\label{v4-open:super-yangian-chain-level}
Construct a Vol IV realization pair for every statement in Vol
I concerning the chain-level chiral coproduct of super-Yangians
\(Y(\mathfrak{sl}(m|n))^{\mathrm{ch}}\), outside the domain of FM230's
absorption into parity-balance super-Whitehead. A concrete super-
Feigin--Frenkel screening construction for \(Y(\mathfrak{sl}(2|1))^{\mathrm{ch}}\)
at generic parameters suffices as the seed.
\end{openproblem}

\begin{openproblem}[Five of six CY-C bridges]
\label{v4-open:cy-c-bridges}
Prove each of the five conditional CY-C constructions \(R_2\)--\(R_6\) of
Vol III (Borcherds multiplicative lift, Mukai-lattice VOA, Kummer
construction, \(\mathcal{N}=(2,2)\) sigma-model half-twist, BLLPR
Schur-limit construction) by proving
the specific convergence isomorphism attaching each construction to the
common Stage-1 \((\Sigma_2,C)\)-specialisation of \(\Phi_3\). Each
discharge is a standalone result: proof of the Igusa cusp form
product convergence against the \(R_1\) normalisation; identification
of the Mukai-lattice OPE with the framed Stage-1 chiral output;
identification of the Kummer-orbifold chiral data with the smooth
\(K3\) chiral specialisation; and the analogous equalities for the
sigma-model and Schur-limit constructions.
\end{openproblem}

\section{Structural Consequences Inherited from Vols I--III}
\label{v4-sec:architectural-sharpenings-inherited}

The realization functor acts on the Vol~I--III theorems in their
stated form. The consequences below specify the independent
comparisons attached to inherited constructions.

\subsection{The four-conductor split and the Koszul-self-dual trinity}
\label{v4-subsec:four-conductor-trinity}

Off the Koszul-self-dual locus the single Koszul conductor
\(K(\mathcal{A}) = -c_{\mathrm{ghost}}(\mathrm{BRST}(\mathcal{A}))\) of Vol~I
equation~(G1) separates into four distinct conductors: \(K\) (BRST
ghost charge), \(K^{\kappa} := \kappa+\kappa^{!}\) (Theorem-C
complementarity ceiling, with values
\((0,13,250/3,25/3,8)\) on
\((\mathsf G,\mathsf L,\mathsf C,\mathsf M,\mathsf B)\), the
\(\mathsf M\)-row value \(25/3\) conditional), \(K_c\)
(central-charge trace), and \(K_g\) (ghost-number anomaly). The trinity
\(K = K_c = K_g\) holds only on the Koszul-self-dual subcategory
admitting both BRST resolution and a Borcherds-lift input on a
Calabi--Yau target.

\noindent\emph{Independent comparison.} The four conductors are
tested against four independent sources: \(K\) against the BRST ghost
resolution, \(K^\kappa\) against the archetype rows of the landscape
census (Kac--Moody trace-form, Virasoro Zamolodchikov norm,
\(\beta\gamma\) contact coefficient, \(\mathcal{W}_N\) Drinfeld--Sokolov
reduction chain, \(K3\) Mukai pairing via Bruinier 2002 Heegner
reciprocity), \(K_c\) against Kac--Wakimoto 1988 central-charge tables,
and \(K_g\) against ghost-partition primary Euler-character
computations.
Disjointness of the four paths is the realization witness of the
four-conductor \emph{separation}; coincidence of three of them on
Koszul-self-dual witnesses the trinity.

\subsection{Universal chain-homotopy normalisation}
\label{v4-subsec:universal-chain-homotopy}

The Loday--Vallette contracting homotopy on the bar--cobar Quillen
pair, specialised to a chirally Koszul \(\mathcal{A}\), rescales by the
complementarity ceiling: \(h_{\mathcal{A}} = h_{\mathrm{LV}}/(\kappa+\kappa^{!})\).
Per archetype: \(h_{\mathsf G}\) trivial (Heisenberg, trace-form
\(\kappa+\kappa^{!}=0\)); \(h_{\mathsf L}\) Sugawara twist \(1/(2(k+h^\vee))\),
diverging at critical level; \(h_{\mathsf C}\) free-field OPE rescaling;
\(h_{\mathsf M}\) Zamolodchikov \(c(5c+22)/10\), singular at \(c=0\),
\(c=-22/5\) Yang--Lee, \(c=-218/45\) secondary Borel.

\noindent\emph{Independent comparison.} The value \(h_{\mathcal{A}}\)
is obtained from the Loday--Vallette universal homotopy
(Loday--Vallette 2012 Theorem~\(11.3.2\)) divided by the Theorem-C
ceiling, and is compared against per-archetype primary-source
computations (Sugawara
twist literature; free-field OPE contractions; Zamolodchikov
\textit{Infinite additional symmetries in two-dimensional conformal
quantum field theory} 1985 for Virasoro norm).

\subsection{Shadow-tower exact closed forms}
\label{v4-subsec:shadow-tower-closed-forms}

The Borel-radius function on the Virasoro shadow tower is
\(|\omega|^{2}(c) = c^{2}(5c+22)/[4(45c+218)]\), large-\(c\) expansion
\(|\omega|(c) = c/6 - 1/27 + O(c^{-1})\), with Stokes pole
\(c_S = -218/45\) (zero of \(180c+872\)). The unweighted shadow coefficient
is the Vol I closed form
\[
S_6(\mathrm{Vir}_c) \;=\; \frac{4(240c+1031)}{c^{3}(5c+22)^{2}}
\]
is forced by the master relation
\(2S_2 S_6 + 2S_3 S_5 + S_4^{2} = 0\) reading off the
\([t^{6}]\) coefficient of \(H^{2}\); the Riccati-metric weighted
companion is
\(\widehat{S}_6^{\mathrm{wt}}(\mathrm{Vir}_c) = 80(45c+193)/[3c^{3}(5c+22)^{2}]\).
The master equation \(H^2 = t^4 Q_c\) is
quadratic-algebraic; the Riccati structure lives on
\(U = (\log H)' - 2/t\), satisfying \(U' + 2U^2 = Q_c''/(2Q_c)\).

\noindent\emph{Independent comparison.} The value
\(|\omega|^2(c)\) is obtained by iterated Borel summation of the shadow
tower on the Riccati-in-\(U\) substitution and is checked at
\(c \in \{1, -22/5, 0, -218/45\}\) against the Borel
radius computed from the Zamolodchikov norm formula
\(c(5c+22)/10\) and the degeneration of the quintic resolvent at
\(c=-22/5\). The Riccati substitution and the Zamolodchikov norm table
are disjoint primary sources.

\subsection{Shadow-Galois Chenevier determinant at \texorpdfstring{\(p=691\)}{p=691}}
\label{v4-subsec:shadow-galois-chenevier-691}

Shadow tower level \(r=11\), Kubota--Leopoldt irregular prime \(p=691\):
\[
\det\rho_{S_{11}(V_1(\mathfrak g)),\,691}
 \;=\; \chi_{\mathrm{cyc}}^{11}\cdot\varepsilon_{B_{12}} \pmod{691},
\]
via Kubota--Leopoldt \(p\)-adic \(L\)-function plus Herbrand--Ribet.

\noindent\emph{Independent comparison.} The determinant
\(\det\rho_{S_{11}}\) at \(p=691\) is computed from the shadow-tower
motivic realisation and checked against \(B_{12}=-691/2730\) and the
Herbrand--Ribet
character table from Washington \textit{Introduction to Cyclotomic
Fields} 1997 (primary source disjoint from the
shadow-tower derivation).

\subsection{Drinfeld double global obstruction class}
\label{v4-subsec:drinfeld-double-obstruction}

For \(g \ge 2\) the universal Drinfeld double \(D(\mathcal{A}, \mathcal{A}^!)\) carries
an explicit cohomological obstruction
\[
 \mathrm{obs}^{(1)}_{\mathrm{double}}
 \;\in\; H^2\bigl(\mathfrak{grt}^{\mathrm{ell}},\,
 \mathfrak{sp}(\mathcal{A})\otimes\mathfrak{sp}(\mathcal{A}^!)\bigr),
 \qquad \omega_{\mathrm{ell}} = \wp(\tau)\langle\cdot,\cdot\rangle_{\mathfrak{sp}},
\]
vanishing if and only if \(r_{\max}(\mathcal{A})=2\) (class \(\mathsf G\) only);
classes \(\mathsf L\), \(\mathsf C\), \(\mathsf M\) are \(g \ge 1\) obstructed.

\noindent\emph{Independent comparison.}
\(\mathrm{obs}^{(1)}_{\mathrm{double}}\) is computed on each archetype via the
Enriquez \(\mathrm{GRT}_1^{\mathrm{ell}}(\mathbb{Q})\) construction (Enriquez
2008 \textit{Elliptic associators}) and compared with vanishing at class~\(\mathsf G\) (Heisenberg)
and nonzero obstruction at classes \(\mathsf L,\mathsf C,\mathsf M\)
via a disjoint computation on the elliptic \(2\)-cocycle
\(\wp(\tau)\langle\cdot,\cdot\rangle_{\mathfrak{sp}}\).

\subsection{Eight-form Gritsenko--Cl\'ery spread and cover assignment}
\label{v4-subsec:eight-form-cover-assignment}

Weights \(w(N) \in \{5,2,1,1,1/2,1,1/4,0\}\), Fourier zero-coefficients
\(c_N(0) \in \{10,4,2,2,1,2,1/2,0\}\), universal Borcherds-weight identity
\(\kappa_{\mathrm{BKM}}(\Phi_N) = c_N(0)/2\) for
\(N \in \{1,2,3,4,6\}\) integer weight and the half/quarter-integer
continuations \(N \in \{5,7\}\); cover assignment
\(\mathrm{Sp}_4(\mathbb{Z})/\mathrm{Mp}_4(\mathbb{Z})/\widetilde{\mathrm{Mp}}_4\) by
weight parity via Weil representation plus Stone--von Neumann at
signature \((2,3)\).

\noindent\emph{Independent comparison.}
\(\kappa_{\mathrm{BKM}}(\Phi_N) = c_N(0)/2\) is evaluated at each row using
\textbf{distinct primary sources}: Gritsenko 1999 \textit{Reflective
modular forms in algebraic geometry} for \(N=1\) (\(\Delta_5\) weight);
Gritsenko--Nikulin 1998 \textit{Automorphic forms and
Lorentzian--Kac--Moody algebras} for \(N \in \{2,3,4,6\}\);
Gritsenko 1994 \(\Gamma_0(5)^+\) for \(N=5\) Borcea--Voisin; Scheithauer
2015 Duke for \(N=7\) quarter-integer; Mongardi--Tari--Wandel 2020 for
\(N=8\) terminal fibre. The eight sources are pairwise distinct,
each serving exactly one row. The identity
\(\kappa_{\mathrm{BKM}} = \kappa_{\mathrm{ch}} + \chi(\mathcal{O}_{\mathrm{fiber}})\)
is not a theorem at any \(N\); the independent comparison rejects
any proof which uses this additive split.

\subsection{Three-factor universal trace identity with \texorpdfstring{\(\zeta\)}{zeta}-regularisation}
\label{v4-subsec:three-factor-UTI}

On the Koszul-self-dual subcategory admitting BRST resolution and
Calabi--Yau target supporting a Borcherds product,
\[
 \mathrm{tr}^{\zeta}_{\mathrm{ghost}}(Q_{\mathrm{BRST}}^2)
 \;=\; \mathrm{tr}_{\mathrm{Pentagon}}
 \;=\; \omega_{\mathrm{Borcherds}}
 \;=\; c_N(0)/2.
\]
\(\zeta\)-regularisation via Dabholkar--Harvey BPS counting (IIA
\(K3\times T^2\), \(\mathcal{N}=4\), \(D=4\) BPS supermultiplicity).
At \(N=1\) four paths converge on \(c_1(0)/2=5\); at \(N=2\) three paths
converge on \(c_2(0)/2=2\).

\noindent\emph{Independent comparison.} Four disjoint calculations
give the four paths at \(N=1\): Vol~I BRST ghost anomaly (\(K(\Phi_3(K3\times E))\));
Vol~III Gritsenko \(\Delta_5\) weight; Vol~II Pentagon face of the
\(\mathrm{SC}^{\mathrm{ch,top}}\) boundary (Heptagon face~4 / Dimofte
slab trace); Dabholkar--Harvey 2007 \textit{BPS states and their
supermultiplicity}. Their sources are pairwise disjoint and their
values all equal \(c_1(0)/2 = 5\), the
common Borcherds--Kac--Moody central-charge value at \(N=1\).

\subsection{Two-stage factorisation of the Calabi--Yau-to-chiral functor}
\label{v4-subsec:two-stage-cy-functor}

\[
 \Phi_d \;=\; \mathrm{Sp}^{\mathrm{ch}}_{\Sigma_{d-1},C}\circ\Phi^{\mathrm{FA}}_d,
\]
with Stage~1 outputting a canonical \(E_d\)-holomorphic factorisation
algebra pinned by Kontsevich--Tamarkin \(E_d\)-formality
(category-level closed at \(d \le 3\), open at \(d \ge 4\)) plus
Costello--Gwilliam--Li locality; Stage~2 specialising via
factorisation homology \(\int_{\Sigma_{d-1}}\) over a \((d-1)\)-cycle,
restriction to a reference curve \(C\). PTVV-forced native operadic
level: \(d=1 \Rightarrow n=\infty\), \(d=2 \Rightarrow n=2\),
\(d \ge 3 \Rightarrow n=1\).

\noindent\emph{Independent comparison.} The Stage-1 calculation builds the
\(E_d\)-FA from the Kontsevich--Tamarkin side (primary source:
Kontsevich 1999 \textit{Operads and motives in deformation
quantization}; Tamarkin 2003 \textit{Action of the Grothendieck--Teichm\"uller
group on the operad of little disks}); the Stage-2 calculation builds the
\(E_1\)-chiral shadow from the factorisation-homology side (primary
source: Ayala--Francis--Tanaka 2017 \textit{Factorization
homology of stratified spaces}; Costello--Gwilliam \textit{Factorization
algebras in quantum field theory} 2021). Pairwise
disjointness is checked, the PTVV shift \((2-d)\) is confirmed via Pantev--To\"en--Vaqui\'e--Vezzosi
2013 primary-source evaluation on each \(d\), and cross-verifies the
Dunn--Lurie native-level dispatch from Lurie \(\mathrm{HA}.5.1.2.2\).

\subsection{K3 Hall--Drinfeld chiral-bialgebra equivalence}
\label{v4-subsec:k3-hall-drinfeld-equivalence}

\[
 \mathrm{Sp}^{\mathrm{ch}}_{K3,E}(\Phi^{\mathrm{FA}}_3(\mathrm{Perf}(K3\times E)))
 \;\simeq\; \mathbf H_{\Delta_5}
\]
is an equivalence at the chiral-bialgebra level, stronger than
graded-dimension coincidence. The six canonical constructions
\(R_1\)--\(R_6\) of Vol~III's six-construction taxonomy
(\(R_1\) framed Stage-1 \(\Phi_3^{(\Sigma_2,C)}\), \(R_2\) Borcherds
multiplicative lift, \(R_3\) Mukai-lattice VOA, \(R_4\) Kummer
construction, \(R_5\) \(\mathcal{N}=(2,2)\) sigma-model half-twist,
\(R_6\) BLLPR Schur limit) form a common limit cone on the chiral
bialgebra \(\mathbf{H}_{\Delta_5}\), with \(R_1\) the proved
generator and \(R_2\)--\(R_6\) carrying convergence hypotheses
(\Cref{v4-sec:cy-c-conditional}).

\noindent\emph{Independent comparison.} The six constructions give
six disjoint calculations of the chiral-bialgebra structure on
\(\Phi_3(D^b(K3\times E))\), and the comparison is natural
isomorphism of the six outputs. The Fake-Monster rank obstruction
excludes any attempt to realise \(\mathrm{II}_{25,1}\)-Leech-lattice chiral data on
a compact CY\(_3\): the bound
\(h^{1,1}(X) \le h^{1,1}(\mathrm{fibre}) + 2h^{0,2}(\mathrm{fibre}) \le 22 < 24\)
is a cohomological obstruction.

\subsection{CoHA-to-\texorpdfstring{\(W_{1+\infty}\)}{W1+inf} triangle with two
\texorpdfstring{\(\lambda\)}{lambda}-parameters}
\label{v4-subsec:coha-triangle-two-lambda}

\(\mathrm{CoHA}(\mathbb{C}^3) = Y^+(\hat{\mathfrak{gl}}_1)\) primitive (Schiffmann--Vasserot
2013); Drinfeld double \(Y(\hat{\mathfrak{gl}}_1)\) affine Yangian of
\(\mathfrak{gl}_1\) in Drinfeld currents (Tsymbaliuk 2017);
\(W_{1+\infty}\) via \(\mathrm{ev}_\lambda\) evaluation (Gaiotto--Rap\v c\'ak
2017). Two \(\lambda\)-parameters are distinguished at all times:
\(\lambda_{\mathrm{Tr}} = (\varepsilon_1+\varepsilon_2)/\varepsilon_3\)
(truncation, fixed to \(-1\) under the CY\(_3\) slice) versus \(\lambda_W\)
(Gaiotto--Rap\v c\'ak central-charge, free). \(\mathrm{CoHA}(\mathbb{C}^3)\)
and \(W_{1+\infty}\) are not equal as algebras.

\noindent\emph{Independent comparison.} Three disjoint character
calculations occur:
CoHA Euler characteristic (Kontsevich--Soibelman 2008); Yangian PBW
character (Drinfeld 1987); \(W_{1+\infty}\) vacuum module character
(Frenkel--Kac--Radul--Wang 1995 \textit{\(W_{1+\infty}\) and
\(W(\mathfrak{gl}_N)\) with central charge \(N\)}). The
\(\lambda\)-parameter discipline is part of the statement:
\(\lambda_{\mathrm{Tr}}\) and \(\lambda_W\) are distinct parameters.

\subsection{Three-tier topologisation discipline}
\label{v4-subsec:three-tier-topologisation}

Tier~I cohomological and Tier~II chain-level on quasi-isomorphic
models are unconditional. Tier~III strict chain-level on the original
complex is closed for classes \(\mathsf G, \mathsf L, \mathsf C\) and
open for class \(\mathsf M\) (the F3 frontier). Sugawara explicit at
\(V_1(\mathfrak{sl}_2)\):
\(T_{\mathrm{Sug}} = (1/6)(\mathord{:}J^0J^0\mathord{:}
+ \mathord{:}J^+J^-\mathord{:} + \mathord{:}J^-J^+\mathord{:})\),
\(G_{\mathrm{Sug}} = (1/3)\sum_a \mathord{:}\bar c_a J^a \mathord{:}\),
\([Q_{\mathrm{tot}}, G_{\mathrm{Sug}}] = T_{\mathrm{Sug}}\) on BRST
cohomology.

\noindent\emph{Independent comparison.} Three comparisons per class, one
per tier: Tier~I (Feigin--Frenkel 2011 cohomological \(E_3\) on
\(H^\ast(\mathrm{Obs}_{\mathrm{hCS}}^{\mathrm{BV}})\)); Tier~II (Fresse 2017
Theorem~\(12.3.\mathrm{A}\) homotopy transfer); Tier~III (explicit
chain-level antighost primitive per archetype). The three tiers have
pairwise disjoint source sets; the class-\(\mathsf M\) Tier~III form is
conditional on F3
resolution.

\subsection{Khan--Zeng finite-jet BV/PVA sector}
\label{v4-subsec:khan-zeng-finite-jet-sector}

For finite-type freely generated nonlinear finite-jet Poisson vertex
algebras (PVAs), Vol~II
constructs the shifted-cotangent BV action on the half-space and
identifies the classical master equation with the PVA Hamiltonian
condition. This is a constructed classical statement. The all-loop
quantum master equation, reflected weight identity, boundary vertex
algebra, and \(E_3\)-lift are conditional on the analytic strong
deformation retract (SDR) datum: finite-jet quotient compatibility,
image-choice Stokes cancellation, field-parity normalisation, and
self-image renormalisation.

\noindent\emph{Independent comparison.} The classical comparison builds
the shifted-cotangent BV action from the PVA Hamiltonian density and
checks CME equality against the lambda-bracket Hamiltonian condition.
The quantum comparison is absent until it supplies
the analytic SDR on every jet-weight quotient, proves compatibility
in the inverse limit, and verifies the reflected weight identity
independently of the classical CME derivation. Thus there are
two forms: constructed classical realization for CME/PVA, and
conditional quantum realization for QME/boundary-VA/\(E_3\).

\subsection{Joyce--KPS compact Calabi--Yau threefold sector}
\label{v4-subsec:joyce-kps-compact-cy3}

For a generic smooth projective Calabi--Yau threefold \(X\), Vol~III
uses Joyce's vertex algebra \(V_X\) on the derived moduli stack and
the Kontsevich--Pandharipande--Soibelman orientation to place MNOP in
the \(E_2\)-centre of the \(E_3\) factorisation algebra on the
constructed module surfaces. The \(S^3\)-framing statement for
generic compact Calabi--Yau Lagrangians is chain-level on the
Costello--Paquette defect-chiral model. A comparison with the global
Hall--Borcherds or Drinfeld envelope is an additional bridge, not a
source of the local construction.

\noindent\emph{Independent comparison.} The compact-CY3 comparison
checks the Joyce \(V_X\) construction and KPS orientation data
against the MNOP centre automorphism. A separate chain-level comparison
checks the \(S^3\)-framing through the defect-chiral model. A third
comparison may test Hall--Borcherds/Drinfeld compatibility, but that
bridge is not a proof source for the Joyce--KPS construction itself.

\subsection{Post-Tate Hadamard support and the \texorpdfstring{\(M_t\)}{Mt}-graph dual}
\label{v4-subsec:post-tate-hadamard-meyer}

The arithmetic spectral-flow branch separates the critical-line
ordinate support
\[
 \Gamma_\xi^{\mathrm{crit}}
 = \{\gamma\in \mathbb R:\xi(1/2+i\gamma)=0\}
\]
from the full Hadamard divisor. Multiplicities are forgotten in this
support. The ordinary \(M_t\)-graph topology (the Hilbert-space
topology induced by the graph norm of the multiplication operator
\(M_t\) on \(L^2(\mathbb R, dt)\), i.e.\ the \((1+t^2)\,dt\) inner
product) does not admit a non-zero continuous graph-dual functional
supported on \(\Gamma_\xi^{\mathrm{crit}}\), even after finite-jet
residue data are allowed. The obstruction is functional analytic:
the graph dual is represented through the resolvent chart by
\(L^2\)-data, while a locally finite critical-line support has
measure zero.

\noindent\emph{Independent comparison.} The ordinary \(M_t\)-graph
dual after the post-Tate quotient gives the \(L^2\)-support
obstruction. A surviving construction must produce a continuous Hadamard
trace-and-jet map
\(\Phi_\xi\to \ell^2(\Gamma_\xi^{\mathrm{crit}},w)\)
via a rigged, Paley--Wiener, de Branges, reproducing-kernel,
nonlocal \(L^2\), or BV--BFV boundary topology, prove post-Tate
compatibility and adjoint-domain compatibility, and verify that the
closed condition arises from the analytic structure and is not
imposed by fiat.

\subsection{Boundary frontiers F1--F10}
\label{v4-subsec:boundary-frontiers}

The ten boundary frontiers are the complement of the present
independent domain: (F1) logarithmic \(\mathcal{W}(p)\) triplet tempering;
(F2) periodic CDG for non-simply-laced admissible levels;
(F3) original-complex Tier~III for class~\(\mathsf M\);
(F4) super-complementarity canonical pairing;
(F5) irregular-singular KZB at \(\overline{\mathcal{M}}_{g,n}\) boundaries;
(F6) Theorem~B at nodal boundary without Koszul hypothesis;
(F7) MC4 completion / H3 failure regime;
(F8) Stage-1 chain level at \(d \ge 4\);
(F9) category-level \(E_d\)-formality at \(d \ge 4\);
(F10) bracket-level BPS Lie algebra on \(K3\times E\).

When any of F1--F10 is proved in Vol~I--III, its realization pair
\(\mathrm{Real}(X, T) = (d, e)\) is the corresponding independent
comparison in Vol~IV. Until then it remains a boundary stratum.
